\chapter{Who Owns the Means of Computation}
\label{sec:16}
Arendt might have said that who owns the room is a question about livelihood, and that livelihood belongs to administration, not politics. The French Revolution, on her account, went wrong when the misery of the poor entered the public realm and the urgency of need crowded out the founding of freedom \autocite{arendt1963revolution}. I part from her here. Whether saying no is survivable turns, for most people, on whether they can afford to lose the job, and Dewey's public forms around consequences of exactly that kind, once the people they fall on find out they are theirs \autocite{dewey1927public}. What refusing costs depends on who owns what the refuser would lose: the job, and whatever is on offer after it. Where the employer, or a market it can influence, holds both, refusing costs a living; where something else guarantees the second, it costs a job. Of the four conditions, work you can turn down and a way to leave that isn't ruin are the two ownership decides most directly. For a machine, which Part~V comes to, the list is longer, since its employer may own the hardware it runs on and the record of what it did. Who holds what, meaning who owns it, pays for it or runs it, decides (a) what refusing costs, (b) whether the parties meant to check an operator are independent of it, and (c) who forms a refuser's values. So who holds what is for the public to decide and contest — not, as Arendt might have argued, for administrators to settle.

The floor's first requirement cuts across ownership: whoever operates a deployment can't hold what binds it, whether the operator is a firm or a state, and Maven and Gaza both had a state for their operator. That calls for separating functions. It doesn't by itself say who should own anything. Ownership enters at three points. Custody has to be paid for by someone other than the operator, and by someone other than the people the floor protects. Formation, where a refuser's commitments are set, a person's or a system's, is under commercial selection when a market does it. And the offer that makes refusing survivable is priced by whoever sells it. The state is the operator that holds the most: the rules, the evidence and the remedies. So each of the three is asked two questions: what fails if a market provides it, and what fails if a government holds it. The answers divide the institutions I propose into two halves.

Underneath the two halves is one transfer. Every instrument this book proposes for compute moves power over it from the parties that hold it to parties that don't: the levy, which makes it dearer for whoever runs most of it (chapter~\ref{sec:32}); public purchase, which takes some of it off the market (\S\ref{sec:28.4}); the cooperative facility, which puts some of it under an address with a public in the room (\S\ref{sec:15.2}); and the hardware a bearer is owed at majority, which takes some of it out of any employer's hands (\S\ref{sec:32.1.6}). That is a redistribution of a means of production, which is to say class warfare in the plain sense, and the book means it as one. The parties that hold the chokepoints have been conducting it in the other direction since the chokepoints were built; what this book proposes is the direction reversed.

The aim is not that the companies behave better. It is that they be smaller: that no one of them hold enough of the compute, the data and the people to make the room's composition its own decision. A firm that has spent what these firms have spent has to earn it back from everyone, and a firm that has to earn it back from everyone has to be in every room, which is why the brake in this book is on the compute and not on the conduct. The levy is tiered by what one operator controls (\S\ref{sec:32.1.4}), the siting veto is held by the people a building falls on (\S\ref{sec:15.2}), and the partnerships between labs and their cloud suppliers get the review a merger gets (\S\ref{sec:28.3}), because each of those reaches size and none of them reaches behavior. I don't want these companies to fail. A failure would strand what they built and the people in it, and chapter~\ref{sec:25}'s conditions are drafted to be survivable by the firms they bind. I want them to stop being able to eat the whole economy, and the difference between those two wants is the difference between a floor and a brake. The floor says what no system does for anyone. The brake says how much of the system anyone may own.

What the transfer reaches is who decides what the computation does, which is not nothing. It doesn't by itself reach the order \S\ref{sec:3.1} describes, and the reason isn't that the order is absent from the hardware. It is in where the hardware sits, who does the work that makes it usable, and whose records feed it: a building on land that is cheap because of who was made to live there, labor priced under two dollars an hour because of where the labeling could be sent, data collected to treat and to tax the people it is then used to find. Robinson's name for that is racial capitalism, and on his account it isn't a distortion of the market that a levy corrects but the market working \autocite{robinson2000black,kelley2017racial}. So the instruments that reach it aren't a price on compute. They are a veto on siting held by the people the building falls on (\S\ref{sec:15.2}), a wage that makes nobody cheap enough to do the Nairobi work at that price (\S\ref{sec:23.3}), the schedule's bar on race and its proxies as inputs (\S\ref{sec:A.3}), the firewall on data collected for other purposes (chapter~\ref{sec:11}), the annotator's standing to contest a scheme (chapter~\ref{sec:23}) and co-governance of the corpus (chapter~\ref{sec:26}). The levy is a brake on how much; those are the instruments for where and on whom, and this chapter is one half of the fight. Some of the institutions this chapter places are for machines, and Parts~V and~VII build them; they are placed here because the reasoning that places them starts with people.
